\section{Implementation}
\label{sec:implementation}

The implementation is a Python 3.12+ package of approximately $16\,500$
lines organised into ten modules and covered by $416$ passing tests
(Section~\ref{sec:impl:testing}). This
section describes the parts of the codebase that a reader would need to see
in order to understand or reproduce the paper's claims.

\subsection{Repository Structure}
\label{sec:impl:repository}

\begin{description}
    \item[\texttt{sfdb.common}] Shared data types (\texttt{SemanticFact},
          \texttt{Entity}, \texttt{Relation}, \texttt{Context}) and the
          engine ABC that both storage engines implement.
    \item[\texttt{sfdb.canonical}] The canonical intermediate
          representation and the bidirectional mapping between it and each
          engine's native format.
    \item[\texttt{sfdb.kg}] The Knowledge Graph engine: dictionary
          encoding, SPO/POS/OSP indexes, reification, and a SPARQL-inspired
          query executor.
    \item[\texttt{sfdb.sheaf}] The sheaf-inspired engine: incremental
          topology, open-set index, stalk index, restriction cache,
          neighbourhood index, global-section constructor, and consistency
          checker.
    \item[\texttt{sfdb.query}] The shared query language, parser, and
          logical plan representation.
    \item[\texttt{sfdb.optimizer}] The shared cost-based optimiser
          (predicate pushdown, dead-operator elimination) with
          engine-specific cost models.
    \item[\texttt{sfdb.datasets}] Synthetic dataset generation and a
          LUBM~\cite{guo2005lubm} loader.
    \item[\texttt{sfdb.benchmark}] The benchmark harness: runners,
          per-engine adapters, metric collectors, and the cross-engine
          verifier.
    \item[\texttt{sfdb.storage}] JSON, MessagePack, and Parquet
          serialisation for facts and results.
    \item[\texttt{sfdb.visualization}] Plot generation used by
          \texttt{scripts/generate\_paper.py} to produce the figures in
          Section~\ref{sec:evaluation}.
\end{description}

\subsection{Mapping to the Formal Model}
\label{sec:mapping:kg-real}
\label{sec:mapping:sh-real}

The KG mapping $\phi_{\mathrm{KG}}$ reifies a fact into a set of triples
against a fresh event identifier. Given a fact of arity $k$, the mapping
emits $k + 4$ triples: a type triple $(e, \texttt{rdf:type}, \texttt{sfdb:Fact})$,
a subject triple, a relation triple, one triple per object slot, a context
triple, and a confidence triple. Provenance and temporal envelopes add
additional triples when present. The inverse mapping $\psi_{\mathrm{KG}}$
gathers all triples sharing the same event identifier and reconstitutes
the original fact.

The sheaf mapping $\phi_{\mathrm{Sh}}$ stores the fact as a single
\texttt{LocalSection} attached to its context. The section retains the
complete tuple of arguments. No decomposition takes place; retrieval by
fact identifier is a single dictionary lookup. The inverse mapping is the
identity on the stored payload.

Both mappings are lossless in the sense of Definition~\ref{def:lossless}.
The round-trip $\psi \circ \phi = \mathrm{id}_\mathcal{F}$ is enforced by
the test suite for every fact inserted during a benchmark run
(\texttt{tests/test\_roundtrip.py}).

\subsection{Configuration and Reproducibility}
\label{sec:impl:config}

The system reads configuration from three sources, in order of increasing
priority: a YAML file (\texttt{configs/benchmark.json} for the default
benchmark suite), environment variables prefixed with \texttt{SFDB\_}, and
programmatic overrides through a \texttt{Config} dataclass. Every
benchmark run seeds Python's PRNG and NumPy's PRNG from the same integer
(default $42$), records the git commit hash, the Python version, the
\texttt{uv.lock} hash, and a checksum of the generated dataset, and stores
them in a per-run reproducibility manifest.

\subsection{Testing}
\label{sec:impl:testing}

The test suite uses \texttt{pytest}. Module tests validate individual
components: topology construction, restriction composition, stalk indexing,
reification round-trips, query planning, and canonicalisation. Cross-engine
tests (\texttt{tests/test\_cross\_engine.py}) execute the same fact stream
and the same query set on both engines and assert canonical equality of the
result sets. Round-trip tests (\texttt{tests/test\_roundtrip.py}) assert
that $\psi \circ \phi = \mathrm{id}$ on both engines for a large synthetic
population. As of the current revision, all $416$ tests pass: the $402$
that require no external service complete in roughly four seconds on a
modern laptop, and a further $14$, in
\texttt{tests/benchmark/}, exercise the live
Jena, Virtuoso, and GraphDB adapters of Section~\ref{sec:eval:jena},
Section~\ref{sec:eval:virtuoso}, and Section~\ref{sec:eval:graphdb} when those external services are
reachable, and are skipped rather than failed otherwise, so their
runtime depends on the responsiveness of servers this artifact does not
control.

\subsection{Engine Adapters for External Systems}
\label{sec:impl:adapters}

The benchmark harness supports pluggable engine adapters. Three adapters
are fully implemented and used in the core verified evaluation:
\texttt{KGEngineAdapter}, \texttt{KGMemEngineAdapter} (the storage-layer
control of Section~\ref{sec:eval:control}, selected by a
\texttt{storage="memory"} engine configuration), and
\texttt{SheafEngineAdapter}. A fourth, \texttt{RdflibEngineAdapter},
wraps the \texttt{rdflib} Python library~\cite{rdflib} and is driven by
SPARQL text; it powers the external reference point of
Section~\ref{sec:eval:rdflib}. An earlier revision named this class
\texttt{JenaEngineAdapter}, which overstated what it wraps --- rdflib is
not Apache Jena --- and we renamed it rather than leave the misleading
name in the artifact.

A fifth, \texttt{JenaTDB2EngineAdapter}, wraps a real Apache Jena TDB2
instance: it speaks the standard SPARQL~1.1 HTTP protocol (SPARQL Update
for inserts, SPARQL Query for reads) to a Fuseki server backing a
disk-based TDB2 dataset, using the identical reification predicates and
SPARQL query text as \texttt{RdflibEngineAdapter} so the two external
comparisons are directly comparable. It powers
Section~\ref{sec:eval:jena}'s production-store reference point. Unlike
every other adapter in this list, the engine behind it is neither
written nor controlled by this project's authors. Scaffolds for
Blazegraph and Neo4j remain and still raise
\texttt{NotImplementedError} on construction; they are place-holders for
planned integration work.
